\documentclass[11pt]{article}
\usepackage[T1]{fontenc}
\usepackage{baskervillef}
\usepackage[varqu,varl,var0]{inconsolata}
\usepackage[scale=.95,type1]{cabin}
\usepackage[baskerville,vvarbb]{newtxmath}
\usepackage[cal=boondoxo]{mathalfa}


\usepackage{hyperref}

\newcommand{\fullcoverage}{\textit{Full coverage}}
\newcommand{\partialcoverage}{\textit{Partial coverage}}
\newcommand{\nocoverage}{\textit{Not covered}}

\begin{document}
\begin{center}\Large\bfseries
  Association for Computing Recommendations \\
  and \textit{Theory of Computation} 
\end{center}

This describes what the text covers of the material that is asked for in the the ACM's
\href{https://www.acm.org/binaries/content/assets/education/cs2013_web_final.pdf}{\textit{Curriculum Guidelines for Undergraduate Programs in Computer Science}}(2013).



\section{AL/Basic Automata Computability and Complexity}

\begin{description}
\item[Topics:]
\begin{itemize}
% \item[Core-Tier1]
\item Finite-state machines  \fullcoverage
\item Regular expressions     \fullcoverage
\item The halting problem      \fullcoverage
% \item[Core-Tier2]
\item Context-free grammars  \partialcoverage
\item Introduction to the P and NP classes and the P vs.\ NP problem  \fullcoverage
\item Introduction to the NP-complete class and exemplary NP-complete problems (e.g., SAT, Knapsack)  \fullcoverage
\end{itemize}
\item[Learning Outcomes:]
\begin{enumerate}
% \item[Core-Tier1]
\item Discuss the concept of finite state machines. [Familiarity]  \fullcoverage
\item Design a deterministic finite state machine to accept a specified language. [Usage]  \fullcoverage
\item Generate a regular expression to represent a specified language. [Usage]  \fullcoverage
\item Explain why the halting problem has no algorithmic solution. [Familiarity]  \fullcoverage
% \item[Core-Tier2]
\item Design a context-free grammar to represent a specified language. [Usage]  \partialcoverage
\item Define the classes P and NP. [Familiarity]  \fullcoverage
\item Explain the significance of NP-completeness. [Familiarity]  \fullcoverage
\end{enumerate}
\end{description}



\section{AL/Advanced Computational Complexity}
\begin{description}
\item[Topics:] 
\begin{itemize}
\item  Review of the classes P and NP; introduce P-space and EXP  \fullcoverage
\item Polynomial hierarchy   \partialcoverage
\item NP-completeness (Cook's theorem)   \fullcoverage
\item Classic NP-complete problems  \fullcoverage
\item Reduction Techniques   \fullcoverage
\end{itemize}
\item[Learning Outcomes:]
\begin{enumerate}
\item Define the classes P and NP. 
[Familiarity]  \fullcoverage
\item Define the P-space class and its relation to the EXP class. [Familiarity]
   \partialcoverage
\item Explain the significance of NP-completeness. [Familiarity]   \fullcoverage
\item Provide examples of classic NP-complete problems. [Familiarity]   \fullcoverage
\item Prove that a problem is NP-complete by reducing a classic known NP-complete problem to it. [Usage]   \fullcoverage
\end{enumerate}
\end{description}




\section{AL/Advanced Automata Theory and Computability}
\begin{description}
\item[Topics:]
\begin{itemize}
\item Sets and languages
  \begin{itemize}
    \item Regular languages   \fullcoverage
    \item Review of deterministic finite automata (DFAs)   \fullcoverage
    \item Nondeterministic finite automata (NFAs)    \fullcoverage
    \item Equivalence of DFAs and NFAs    \fullcoverage
    \item Review of regular expressions; their equivalence to finite automata  \fullcoverage
    \item Closure properties   \fullcoverage
    \item Proving languages non-regular, via the pumping lemma or alternative means   \fullcoverage
  \end{itemize}
\item Context-free languages
  \begin{itemize}
    \item Push-down automata (PDAs)   \fullcoverage
    \item Relationship of PDAs and context-free grammars  \partialcoverage
    \item Properties of context-free languages   \partialcoverage
  \end{itemize}
\item Turing machines, or an equivalent formal model of universal computation  \fullcoverage
\item Nondeterministic Turing machines   \fullcoverage
\item Chomsky hierarchy   \nocoverage
\item The Church-Turing thesis   \fullcoverage
\item Computability      \fullcoverage
\item Rice's Theorem      \fullcoverage
\item Examples of uncomputable functions    \fullcoverage
\item Implications of uncomputability    \fullcoverage
\end{itemize}
\item[Learning Outcomes:]
\begin{enumerate}
\item Determine a language's place in the Chomsky hierarchy (regular, context-free, recursively enumerable).   \nocoverage
[Assessment]
\item Convert among equivalently powerful notations for a language, including among DFAs, NFAs, and regular
expressions, and between PDAs and CFGs. [Usage]  \partialcoverage
\item Explain the Church-Turing thesis and its significance. [Familiarity]  \fullcoverage
\item Explain Rice's Theorem and its significance. [Familiarity]   \fullcoverage
\item Provide examples of uncomputable functions. [Familiarity]   \fullcoverage
\item Prove that a problem is uncomputable by reducing a classic known uncomputable problem to it. [Usage]   \fullcoverage
\end{enumerate}
\end{description}

\end{document}
